\chapter{Strings}
\label{chap:strings}

The \texttt{std::strings} module provides a comprehensive string handling system designed for Amiga development. It offers both borrowed views (zero-copy) and owned strings (heap-allocated), plus AmigaDOS-specific types like BCPL strings.

\section{Overview}

Novus provides multiple string types for different use cases:

\begin{table}[h]
\centering
\begin{tabular}{llll}
\textbf{Type} & \textbf{Size} & \textbf{Allocation} & \textbf{Use Case} \\
\hline
\texttt{Str} & 8 bytes & None (view) & Parameters, slices \\
\texttt{String} & 12 bytes & Heap (Vec) & Building, ownership \\
\texttt{NameString} & 36 bytes & Stack (32 chars) & Short names \\
\texttt{GadgetString} & 132 bytes & Stack (128 chars) & UI text buffers \\
\texttt{PathString} & 260 bytes & Stack (256 chars) & File paths \\
\texttt{StringBuilder} & 12 bytes & Heap & Efficient concatenation \\
\texttt{BStr} & 4 bytes & Heap (BCPL) & AmigaDOS interop \\
\end{tabular}
\caption{String types and their characteristics}
\end{table}

\begin{notebox}
All Novus strings are Latin-1 (ISO-8859-1) encoded, matching AmigaOS conventions. ASCII is a subset of Latin-1, so ASCII text works naturally. UTF-8 support may be added in future versions.
\end{notebox}

%%%%%%%%%%%%%%%%%%%%%%%%%%%%%%%%%%%%%%%%%%%%%%%%%%%%%%%%%%%%%%%%%%%%%%%%%%%%%%%
\section{Str --- Borrowed String Slice}
\label{sec:str}
%%%%%%%%%%%%%%%%%%%%%%%%%%%%%%%%%%%%%%%%%%%%%%%%%%%%%%%%%%%%%%%%%%%%%%%%%%%%%%%

\texttt{Str} is a ``fat pointer'' containing a pointer and length. It does NOT own the data---it's just a view into existing memory. This makes it extremely efficient for passing strings to functions.

\begin{lstlisting}
from std::strings::core import Str

fn greet(name: Str) {
    // name is just a view - no allocation, no copy
    print_str(name)
}

fn main() {
    // Create Str from C string literal
    match Str::from_cstr("Hello, World!") {
        Option::Some(s) => greet(s),
        Option::None => {}
    }
}
\end{lstlisting}

\subsection{Creating Str Values}

\begin{lstlisting}
// From null-terminated C string (scans for length)
let s1 = Str::from_cstr("Hello")?

// From raw pointer and known length (unsafe - must be valid)
let s2 = Str::from_raw(ptr, 5)
\end{lstlisting}

\begin{warningbox}
\texttt{Str::from\_raw()} does not validate the pointer. The caller must ensure \texttt{ptr} points to valid memory of at least \texttt{len} bytes.
\end{warningbox}

\subsection{Basic Operations}

\begin{table}[h]
\centering
\begin{tabular}{lll}
\textbf{Method} & \textbf{Returns} & \textbf{Description} \\
\hline
\texttt{len()} & \texttt{u32} & Length in bytes \\
\texttt{is\_empty()} & \texttt{bool} & Check if length is zero \\
\texttt{as\_ptr()} & \texttt{*u8} & Get raw pointer \\
\texttt{get(idx)} & \texttt{Option<u8>} & Get byte at index \\
\texttt{as\_cstr()} & \texttt{*u8} & Get C string pointer \\
\end{tabular}
\caption{\texttt{Str} basic methods}
\end{table}

\subsection{Slicing}

Create substrings without copying:

\begin{lstlisting}
let text = Str::from_cstr("Hello, World!")?

// Slice with start and end indices
let hello = text.slice(0, 5)?    // "Hello"
let world = text.slice(7, 12)?   // "World"

// Slice from index to end
let rest = text.slice_from(7)?   // "World!"

// Slice from start to index
let prefix = text.slice_to(5)?   // "Hello"
\end{lstlisting}

\begin{notebox}
Slicing returns \texttt{Result<Str, StringError>} with bounds checking. The slice is a view into the original data---no allocation occurs.
\end{notebox}

\subsection{Comparison}

\begin{lstlisting}
let a = Str::from_cstr("hello")?
let b = Str::from_cstr("HELLO")?
let c = Str::from_cstr("hello")?

// Exact equality
if a.equals(c) {
    // true - same bytes
}

// Case-insensitive equality
if a.eq_ignore_case(b) {
    // true - ignores ASCII case
}

// Lexicographic comparison (returns i32: <0, 0, or >0)
let cmp = a.compare(b)  // negative (lowercase < uppercase in ASCII)
\end{lstlisting}

\begin{comingfromc}
On 68020+, \texttt{equals()} uses optimized inline assembly with \texttt{cmpm.l} for $\sim$4x speedup on strings $\geq$ 8 bytes. The compiler automatically uses byte-by-byte comparison on 68000/68010.
\end{comingfromc}

\subsection{Searching}

\begin{lstlisting}
let text = Str::from_cstr("Hello, World!")?

// Check prefix/suffix
if text.starts_with(Str::from_cstr("Hello")?) {
    // true
}
if text.ends_with(Str::from_cstr("!")?) {
    // true
}

// Case-insensitive prefix check
if text.starts_with_ignore_case(Str::from_cstr("HELLO")?) {
    // true
}

// Find substring
match text.find(Str::from_cstr("World")?) {
    Option::Some(idx) => {},  // idx == 7
    Option::None => {}
}

// Check if contains substring
if text.contains(Str::from_cstr(",")?) {
    // true
}

// Find single byte
match text.find_byte($2C) {  // ','
    Option::Some(idx) => {},  // idx == 5
    Option::None => {}
}

// Find last occurrence of byte
match text.rfind_byte($6F) {  // 'o'
    Option::Some(idx) => {},  // idx == 8 (second 'o')
    Option::None => {}
}
\end{lstlisting}

\subsection{Trimming}

\begin{lstlisting}
let padded = Str::from_cstr("  hello  ")?

let trimmed = padded.trim()        // "hello"
let left = padded.trim_start()     // "hello  "
let right = padded.trim_end()      // "  hello"
\end{lstlisting}

Trimming removes ASCII whitespace: space (\texttt{\$20}), tab (\texttt{\$09}), LF (\texttt{\$0A}), and CR (\texttt{\$0D}).

%%%%%%%%%%%%%%%%%%%%%%%%%%%%%%%%%%%%%%%%%%%%%%%%%%%%%%%%%%%%%%%%%%%%%%%%%%%%%%%
\section{String --- Owned Heap String}
\label{sec:string-type}
%%%%%%%%%%%%%%%%%%%%%%%%%%%%%%%%%%%%%%%%%%%%%%%%%%%%%%%%%%%%%%%%%%%%%%%%%%%%%%%

\texttt{String} is a heap-allocated, growable, mutable string backed by \texttt{Vec<u8>}. Use it when you need to build strings dynamically or own string data.

\begin{lstlisting}
from std::strings::core import String, Str

fn build_greeting(name: Str) -> Result<String, StringError> {
    var greeting = String::new()
    greeting.push_str(Str::from_cstr("Hello, ")?)?
    greeting.push_str(name)?
    greeting.push_byte($21)?  // '!'
    return Result::Ok(greeting)
}
\end{lstlisting}

\subsection{Creating Strings}

\begin{lstlisting}
// Empty string (no allocation until first push)
var s1 = String::new()

// Pre-allocated capacity
var s2 = String::with_capacity(256)?

// From Str (allocates and copies)
let s3 = String::new_from_str(Str::from_cstr("Hello")?)?
\end{lstlisting}

\subsection{Core Operations}

\begin{table}[h]
\centering
\begin{tabular}{lll}
\textbf{Method} & \textbf{Returns} & \textbf{Description} \\
\hline
\texttt{push\_str(s)} & \texttt{Result<(), StringError>} & Append Str \\
\texttt{push\_byte(b)} & \texttt{Result<(), StringError>} & Append single byte \\
\texttt{len()} & \texttt{u32} & Current length \\
\texttt{capacity()} & \texttt{u32} & Current capacity \\
\texttt{is\_empty()} & \texttt{bool} & Check if empty \\
\texttt{clear()} & \texttt{()} & Clear contents (keep capacity) \\
\texttt{reserve(n)} & \texttt{Result<(), StringError>} & Reserve additional bytes \\
\texttt{as\_str()} & \texttt{Str} & Borrow as Str slice \\
\texttt{as\_ptr()} & \texttt{*u8} & Get raw pointer \\
\texttt{as\_mut\_ptr()} & \texttt{*u8} & Get mutable pointer \\
\end{tabular}
\caption{\texttt{String} method summary}
\end{table}

\subsection{Memory Management}

\texttt{String} implements the \texttt{Drop} trait for automatic cleanup:

\begin{lstlisting}
fn process() -> Result<(), StringError> {
    var message = String::with_capacity(100)?
    message.push_str(Str::from_cstr("Processing...")?)?

    // ... use message ...

    return Result::Ok(())
}  // message.drop() called automatically
\end{lstlisting}

\subsection{Converting Between Types}

\begin{lstlisting}
// Str -> String (allocates)
let owned = str_slice.to_string()?

// String -> Str (borrows)
let borrowed = owned.as_str()

// Both implement AsRef<Str>
fn accepts_str<T: AsRef<Str>>(s: T) {
    let slice: Str = s.as_ref()
}
\end{lstlisting}

%%%%%%%%%%%%%%%%%%%%%%%%%%%%%%%%%%%%%%%%%%%%%%%%%%%%%%%%%%%%%%%%%%%%%%%%%%%%%%%
\section{Fixed-Size Strings}
\label{sec:fixed-strings}
%%%%%%%%%%%%%%%%%%%%%%%%%%%%%%%%%%%%%%%%%%%%%%%%%%%%%%%%%%%%%%%%%%%%%%%%%%%%%%%

For embedded systems work, heap allocation can be expensive or unpredictable. Novus provides stack-allocated fixed-size strings for common use cases.

\subsection{NameString (32 bytes)}

For short identifiers, variable names, and labels:

\begin{lstlisting}
from std::strings::core import NameString, Str

var name = NameString::new()
name.push_str(Str::from_cstr("Player1")?)?

if name.len() < name.capacity() {
    name.push_byte($31)?  // '1'
}

let slice: Str = name.as_str()
\end{lstlisting}

\subsection{GadgetString (128 bytes)}

For Intuition gadget text buffers and UI strings:

\begin{lstlisting}
from std::strings::core import GadgetString, Str

var button_text = GadgetString::new()
button_text.push_str(Str::from_cstr("Click Here")?)?

// Use for gadget creation
let label_ptr: *u8 = button_text.as_mut_ptr()
\end{lstlisting}

\subsection{PathString (256 bytes)}

For AmigaDOS file paths (max path length is typically 256 bytes):

\begin{lstlisting}
from std::strings::core import PathString, Str

var path = PathString::new()
path.push_str(Str::from_cstr("Work:Projects/")?)?
path.push_str(Str::from_cstr("MyApp/main.novus")?)?

// Get as Str for DOS functions
let path_slice: Str = path.as_str()
\end{lstlisting}

\begin{tipbox}
Fixed strings return \texttt{Result::Err(StringError::TooLong)} if you exceed their capacity. Always check the result or use \texttt{capacity()} to verify space.
\end{tipbox}

%%%%%%%%%%%%%%%%%%%%%%%%%%%%%%%%%%%%%%%%%%%%%%%%%%%%%%%%%%%%%%%%%%%%%%%%%%%%%%%
\section{StringBuilder --- Efficient Construction}
\label{sec:stringbuilder}
%%%%%%%%%%%%%%%%%%%%%%%%%%%%%%%%%%%%%%%%%%%%%%%%%%%%%%%%%%%%%%%%%%%%%%%%%%%%%%%

\texttt{StringBuilder} provides efficient string construction with built-in formatting methods for numbers:

\begin{lstlisting}
from std::strings::core import StringBuilder, Str

fn format_status(x: i32, y: i32, health: u32) -> Result<String, StringError> {
    var sb = StringBuilder::new()?

    sb.push_str(Str::from_cstr("Position: (")?)?
    sb.push_i32(x)?
    sb.push_str(Str::from_cstr(", ")?)?
    sb.push_i32(y)?
    sb.push_str(Str::from_cstr(") Health: ")?)?
    sb.push_u32(health)?

    return Result::Ok(sb.build())
}
\end{lstlisting}

\subsection{Formatting Methods}

\begin{table}[h]
\centering
\begin{tabular}{lll}
\textbf{Method} & \textbf{Parameter} & \textbf{Description} \\
\hline
\texttt{push\_str(s)} & \texttt{Str} & Append string \\
\texttt{push\_byte(b)} & \texttt{u8} & Append single byte \\
\texttt{push\_char(c)} & \texttt{u8} & Append character (alias) \\
\texttt{push\_i32(n)} & \texttt{i32} & Format signed integer \\
\texttt{push\_u32(n)} & \texttt{u32} & Format unsigned integer \\
\texttt{push\_hex(n)} & \texttt{u32} & Format as lowercase hex \\
\texttt{push\_hex\_upper(n)} & \texttt{u32} & Format as uppercase hex \\
\texttt{push\_bool(b)} & \texttt{bool} & Format as ``true''/``false'' \\
\end{tabular}
\caption{\texttt{StringBuilder} formatting methods}
\end{table}

\begin{lstlisting}
var sb = StringBuilder::new()?

sb.push_str(Str::from_cstr("Address: $")?)?
sb.push_hex_upper(0xDFF000)?  // "DFF000"

sb.push_str(Str::from_cstr(" Active: ")?)?
sb.push_bool(true)?           // "true"

let result: String = sb.build()
// "Address: $DFF000 Active: true"
\end{lstlisting}

\begin{notebox}
\texttt{StringBuilder::new()} pre-allocates 256 bytes. Use \texttt{with\_capacity(n)} if you know you need more or less.
\end{notebox}

\begin{warningbox}
\texttt{StringBuilder::build()} \textbf{consumes} the builder---you cannot use the builder after calling \texttt{build()}. The \texttt{consuming} keyword in the signature enforces this at compile time.
\end{warningbox}

%%%%%%%%%%%%%%%%%%%%%%%%%%%%%%%%%%%%%%%%%%%%%%%%%%%%%%%%%%%%%%%%%%%%%%%%%%%%%%%
\section{BStr --- BCPL Strings for AmigaDOS}
\label{sec:bstr}
%%%%%%%%%%%%%%%%%%%%%%%%%%%%%%%%%%%%%%%%%%%%%%%%%%%%%%%%%%%%%%%%%%%%%%%%%%%%%%%

AmigaDOS uses BCPL-style strings for many operations. These have a length byte prefix and use ``BPTR'' addressing (address divided by 4).

\begin{lstlisting}
from std::strings::core import BStr, Str

fn dos_example() -> Result<(), StringError> {
    // Create BCPL string from Str
    let bstr = BStr::new_from_str(Str::from_cstr("Work:Projects")?)?

    // Get BPTR for DOS functions
    let bptr: u32 = bstr.as_bptr()

    // Get actual memory address
    let ptr: *u8 = bstr.as_ptr()

    // Get length
    let len: u32 = bstr.len()

    // View as Str (excludes length byte)
    let slice: Str = bstr.as_str()

    return Result::Ok(())
}  // bstr automatically freed via Drop
\end{lstlisting}

\begin{warningbox}
BCPL strings have a maximum length of 255 bytes (stored in a single byte). \texttt{BStr::new\_from\_str()} returns \texttt{Err(StringError::TooLong)} if the input exceeds this limit.
\end{warningbox}

\subsection{BCPL String Format}

\begin{lstlisting}
// Memory layout of BStr for "Hello" (5 bytes)
// Address 0x1000: [05][H][e][l][l][o]
//                  ^-- length byte
// BPTR = 0x1000 / 4 = 0x400
\end{lstlisting}

\begin{comingfromc}
In C, you'd manage BPTR conversion manually:
\begin{lstlisting}[language=C]
BPTR bptr = MKBADDR(str);  // str / 4
char *str = BADDR(bptr);   // bptr * 4
\end{lstlisting}
\texttt{BStr} handles this automatically and provides RAII cleanup.
\end{comingfromc}

\begin{notebox}
\texttt{BStr} has a \texttt{free()} method for explicit cleanup, but this is \textbf{deprecated}. Prefer letting the \texttt{Drop} trait handle cleanup automatically when the \texttt{BStr} goes out of scope.
\end{notebox}

%%%%%%%%%%%%%%%%%%%%%%%%%%%%%%%%%%%%%%%%%%%%%%%%%%%%%%%%%%%%%%%%%%%%%%%%%%%%%%%
\section{String Parsing}
\label{sec:string-parsing}
%%%%%%%%%%%%%%%%%%%%%%%%%%%%%%%%%%%%%%%%%%%%%%%%%%%%%%%%%%%%%%%%%%%%%%%%%%%%%%%

\texttt{Str} provides methods to parse common types from string representations:

\subsection{Integer Parsing}

\begin{lstlisting}
let decimal = Str::from_cstr("12345")?
let negative = Str::from_cstr("-42")?
let hex1 = Str::from_cstr("0xFF00")?
let hex2 = Str::from_cstr("ff00")?

// Parse unsigned
match decimal.parse_u32() {
    Result::Ok(n) => {},    // n == 12345
    Result::Err(e) => {}
}

// Parse signed (handles - and +)
match negative.parse_i32() {
    Result::Ok(n) => {},    // n == -42
    Result::Err(e) => {}
}

// Parse hex (case-insensitive, optional 0x prefix)
match hex1.parse_hex() {
    Result::Ok(n) => {},    // n == 0xFF00
    Result::Err(e) => {}
}
\end{lstlisting}

\subsection{Boolean Parsing}

\begin{lstlisting}
// Accepts: "true", "false", "1", "0" (case-insensitive)
let t1 = Str::from_cstr("true")?.parse_bool()?   // true
let t2 = Str::from_cstr("TRUE")?.parse_bool()?   // true
let t3 = Str::from_cstr("1")?.parse_bool()?      // true
let f1 = Str::from_cstr("false")?.parse_bool()?  // false
let f2 = Str::from_cstr("0")?.parse_bool()?      // false
\end{lstlisting}

\subsection{Parse Errors}

\begin{lstlisting}
pub enum ParseError {
    Empty,              // Empty string
    InvalidChar(u32),   // Bad character at index
    Overflow,           // Value too large
    Underflow,          // Value too small (signed)
    InvalidBool,        // Not "true"/"false"/"1"/"0"
}
\end{lstlisting}

%%%%%%%%%%%%%%%%%%%%%%%%%%%%%%%%%%%%%%%%%%%%%%%%%%%%%%%%%%%%%%%%%%%%%%%%%%%%%%%
\section{Splitting and Iteration}
\label{sec:string-iteration}
%%%%%%%%%%%%%%%%%%%%%%%%%%%%%%%%%%%%%%%%%%%%%%%%%%%%%%%%%%%%%%%%%%%%%%%%%%%%%%%

\subsection{Splitting by Delimiter}

\begin{lstlisting}
let csv = Str::from_cstr("apple,banana,cherry")?

var iter = csv.split($2C)  // ',' = 0x2C
while true {
    match iter.next() {
        Option::Some(part) => {
            // part: "apple", "banana", "cherry"
            print_str(part)
        },
        Option::None => break
    }
}
\end{lstlisting}

\subsection{Splitting into Lines}

\begin{lstlisting}
let text = Str::from_cstr("Line 1\nLine 2\r\nLine 3")?

var lines = text.lines()
while true {
    match lines.next() {
        Option::Some(line) => {
            // Handles LF, CR, and CRLF
            print_str(line)
        },
        Option::None => break
    }
}
\end{lstlisting}

\begin{notebox}
\texttt{SplitIter} and \texttt{LinesIter} implement the \texttt{Iterator<Str>} trait, returning \texttt{Str} slices into the original string. No allocation occurs during iteration.
\end{notebox}

%%%%%%%%%%%%%%%%%%%%%%%%%%%%%%%%%%%%%%%%%%%%%%%%%%%%%%%%%%%%%%%%%%%%%%%%%%%%%%%
\section{String Transformations}
\label{sec:string-transform}
%%%%%%%%%%%%%%%%%%%%%%%%%%%%%%%%%%%%%%%%%%%%%%%%%%%%%%%%%%%%%%%%%%%%%%%%%%%%%%%

\subsection{Case Conversion}

\begin{lstlisting}
let text = Str::from_cstr("Hello, World!")?

// Convert to uppercase (allocates new String)
let upper: String = text.to_uppercase()?
// "HELLO, WORLD!"

// Convert to lowercase (allocates new String)
let lower: String = text.to_lowercase()?
// "hello, world!"
\end{lstlisting}

Case conversion handles both ASCII (A-Z, a-z) and Latin-1 extended characters (\`{a}-\"{o}, \`{A}-\"{O}, etc.).

\subsection{Search and Replace}

\begin{lstlisting}
let text = Str::from_cstr("foo bar foo baz foo")?

// Replace all occurrences
let result: String = text.replace(
    Str::from_cstr("foo")?,
    Str::from_cstr("qux")?
)?
// "qux bar qux baz qux"
\end{lstlisting}

\subsection{Joining Strings}

\begin{lstlisting}
from std::strings::core import str_join, Str

// Join array of Str with separator
var parts: [Str; 3] = [
    Str::from_cstr("SYS")?,
    Str::from_cstr("C")?,
    Str::from_cstr("Dir")?
]

let path: String = str_join((*Str)&parts, 3, Str::from_cstr("/")?)?
// "SYS/C/Dir"
\end{lstlisting}

%%%%%%%%%%%%%%%%%%%%%%%%%%%%%%%%%%%%%%%%%%%%%%%%%%%%%%%%%%%%%%%%%%%%%%%%%%%%%%%
\section{AmigaDOS Path Manipulation}
\label{sec:path-manipulation}
%%%%%%%%%%%%%%%%%%%%%%%%%%%%%%%%%%%%%%%%%%%%%%%%%%%%%%%%%%%%%%%%%%%%%%%%%%%%%%%

Novus provides AmigaDOS-aware path utilities:

\begin{lstlisting}
from std::strings::core import (
    path_split, path_filename, path_dirname,
    path_is_absolute, path_join, PathComponents
)

let full_path = Str::from_cstr("Work:Projects/Novus/main.novus")?

// Split into components
let parts: PathComponents = path_split(full_path)
// parts.volume = Some("Work:")
// parts.directory = Some("Projects/Novus")
// parts.filename = Some("main.novus")

// Get just the filename
match path_filename(full_path) {
    Option::Some(name) => {},  // "main.novus"
    Option::None => {}
}

// Get the directory
match path_dirname(full_path) {
    Option::Some(dir) => {},   // "Work:Projects/Novus"
    Option::None => {}
}

// Check if absolute (has volume)
if path_is_absolute(full_path) {
    // true - has "Work:"
}

// Join paths (handles separators correctly)
let joined: String = path_join(
    Str::from_cstr("Work:Projects")?,
    Str::from_cstr("NewApp")?
)?
// "Work:Projects/NewApp"
\end{lstlisting}

\begin{tipbox}
AmigaDOS path conventions:
\begin{itemize}
    \item Volume names end with colon: \texttt{SYS:}, \texttt{Work:}, \texttt{RAM:}
    \item Directories separated by forward slash: \texttt{/}
    \item Parent directory: \texttt{/} (just slash, unlike \texttt{..} in Unix)
    \item Current directory: \texttt{""} (empty string)
\end{itemize}
\end{tipbox}

%%%%%%%%%%%%%%%%%%%%%%%%%%%%%%%%%%%%%%%%%%%%%%%%%%%%%%%%%%%%%%%%%%%%%%%%%%%%%%%
\section{Display and Debug Traits}
\label{sec:display-debug}
%%%%%%%%%%%%%%%%%%%%%%%%%%%%%%%%%%%%%%%%%%%%%%%%%%%%%%%%%%%%%%%%%%%%%%%%%%%%%%%

The \texttt{std::strings::format} module provides traits for formatted output:

\begin{lstlisting}
from std::strings::format import Display, Debug, Formatter

// Display - user-facing output
pub trait Display {
    fn fmt(&self, f: &var Formatter) -> bool
}

// Debug - programmer-facing debug output
pub trait Debug {
    fn fmt(&self, f: &var Formatter) -> bool
}
\end{lstlisting}

\subsection{Implementing Display}

\begin{lstlisting}
struct Point {
    x: i32,
    y: i32,
}

impl Display for Point {
    fn fmt(&self, f: &var Formatter) -> bool {
        if !f.write_char($28) { return false }  // '('
        // ... write x ...
        if !f.write_str(Str::from_cstr(", ")?) { return false }
        // ... write y ...
        if !f.write_char($29) { return false }  // ')'
        return true
    }
}
\end{lstlisting}

\subsection{Primitive Implementations}

All primitive types implement \texttt{Display} and \texttt{Debug}:

\begin{itemize}
    \item \texttt{i8}, \texttt{i16}, \texttt{i32}, \texttt{i64} --- decimal representation
    \item \texttt{u8}, \texttt{u16}, \texttt{u32}, \texttt{u64} --- decimal representation
    \item \texttt{bool} --- ``true'' or ``false''
    \item \texttt{Str}, \texttt{String} --- string contents
\end{itemize}

\subsection{StackFormatter}

For zero-allocation formatting (e.g., in interrupt handlers), use \texttt{StackFormatter}:

\begin{lstlisting}
from std::strings::format import StackFormatter

fn format_on_stack() {
    var fmt = StackFormatter::new()  // 256-byte stack buffer

    fmt.write_str(Str::from_cstr("Value: ")?)
    // ... format values ...

    let result: Str = fmt.as_str()
    // WARNING: result only valid while fmt is in scope!
}
\end{lstlisting}

\begin{warningbox}
\texttt{StackFormatter::as\_str()} returns a \texttt{Str} pointing into the formatter's stack buffer. This reference becomes invalid when the formatter goes out of scope.
\end{warningbox}

%%%%%%%%%%%%%%%%%%%%%%%%%%%%%%%%%%%%%%%%%%%%%%%%%%%%%%%%%%%%%%%%%%%%%%%%%%%%%%%
\section{Error Types}
\label{sec:string-errors}
%%%%%%%%%%%%%%%%%%%%%%%%%%%%%%%%%%%%%%%%%%%%%%%%%%%%%%%%%%%%%%%%%%%%%%%%%%%%%%%

\subsection{StringError}

\begin{lstlisting}
pub enum StringError {
    OutOfBounds(u32, u32),  // (index, length)
    InvalidUtf8(u32),       // position
    AllocationFailed,
    TooLong(u32),           // length (for BStr: max 255)
    InvalidFormat,
    NullPointer,
}
\end{lstlisting}

\subsection{ParseError}

\begin{lstlisting}
pub enum ParseError {
    Empty,              // Empty string
    InvalidChar(u32),   // position of bad character
    Overflow,           // Value too large
    Underflow,          // Value too small (signed)
    InvalidBool,        // Not true/false/1/0
}
\end{lstlisting}

%%%%%%%%%%%%%%%%%%%%%%%%%%%%%%%%%%%%%%%%%%%%%%%%%%%%%%%%%%%%%%%%%%%%%%%%%%%%%%%
\section{C FFI Functions}
\label{sec:string-ffi}
%%%%%%%%%%%%%%%%%%%%%%%%%%%%%%%%%%%%%%%%%%%%%%%%%%%%%%%%%%%%%%%%%%%%%%%%%%%%%%%

For low-level operations and C library interop, standard C string functions are available:

\begin{lstlisting}
extern fn strlen(str: *u8) -> u32
extern fn strcmp(s1: *u8, s2: *u8) -> i32
extern fn strncmp(s1: *u8, s2: *u8, n: u32) -> i32
extern fn strcpy(dest: *u8, src: *u8) -> *u8
extern fn strncpy(dest: *u8, src: *u8, n: u32) -> *u8
extern fn strcat(dest: *u8, src: *u8) -> *u8
extern fn strncat(dest: *u8, src: *u8, n: u32) -> *u8
extern fn strchr(str: *u8, ch: i32) -> *u8
extern fn strrchr(str: *u8, ch: i32) -> *u8
extern fn strstr(haystack: *u8, needle: *u8) -> *u8
extern fn memmove(dest: *u8, src: *u8, n: u32) -> *u8
extern fn memset(ptr: *u8, value: i32, n: u32) -> *u8
extern fn memcmp(s1: *u8, s2: *u8, n: u32) -> i32
\end{lstlisting}

\begin{warningbox}
These functions provide no bounds checking and can easily cause buffer overflows. Prefer the safe Novus string types whenever possible. These are provided for performance-critical code and C library interop.
\end{warningbox}

%%%%%%%%%%%%%%%%%%%%%%%%%%%%%%%%%%%%%%%%%%%%%%%%%%%%%%%%%%%%%%%%%%%%%%%%%%%%%%%
\section{Trait Implementations}
\label{sec:string-traits}
%%%%%%%%%%%%%%%%%%%%%%%%%%%%%%%%%%%%%%%%%%%%%%%%%%%%%%%%%%%%%%%%%%%%%%%%%%%%%%%

String types implement several important traits for idiomatic usage:

\subsection{Clone}

\texttt{String} and \texttt{StringBuilder} implement \texttt{Clone} for deep copying:

\begin{lstlisting}
let original = String::new_from_str(Str::from_cstr("Hello")?)?

// Clone creates a complete copy (allocates new memory)
match original.clone() {
    Option::Some(copy) => {
        // copy is independent of original
    },
    Option::None => {
        // allocation failed
    }
}
\end{lstlisting}

\subsection{Write}

\texttt{String}, \texttt{StringBuilder}, and \texttt{PathString} implement the \texttt{Write} trait for generic output:

\begin{lstlisting}
from std::core import Write

// Generic function that writes to any Write implementor
fn write_header<W: Write>(w: &var W) -> bool {
    return w.write_bytes("Header: ", 8)
}

var s = String::new()
write_header(&s)  // Works with String

var sb = StringBuilder::new()?
write_header(&sb)  // Works with StringBuilder

var path = PathString::new()
write_header(&path)  // Works with PathString
\end{lstlisting}

\subsection{Ord (Ordering)}

\texttt{Str} implements \texttt{Ord} for lexicographic comparison:

\begin{lstlisting}
let a = Str::from_cstr("apple")?
let b = Str::from_cstr("banana")?

// Compare using Ord trait
let cmp: i32 = a.cmp(&b)
if cmp < 0 {
    // a comes before b lexicographically
}
\end{lstlisting}

\subsection{Iterator}

\texttt{SplitIter} and \texttt{LinesIter} implement \texttt{Iterator<Str>}:

\begin{lstlisting}
let csv = Str::from_cstr("a,b,c")?

// Iterators implement Iterator<Str> trait
var iter = csv.split($2C)  // ','

// Can use in while loops with next()
while true {
    match iter.next() {
        Option::Some(part) => process(part),
        Option::None => break
    }
}
\end{lstlisting}

\subsection{AsRef}

Multiple types implement \texttt{AsRef<Str>} for unified string access:

\begin{lstlisting}
// These all implement AsRef<Str>:
// - Str (identity)
// - String
// - StringBuilder
// - PathString

fn print_anything<T: AsRef<Str>>(s: T) {
    let slice: Str = s.as_ref()
    // Now use slice...
}

let owned = String::new_from_str(Str::from_cstr("Hello")?)?
print_anything(owned)  // Works!

let path = PathString::new()
print_anything(path)   // Also works!
\end{lstlisting}

%%%%%%%%%%%%%%%%%%%%%%%%%%%%%%%%%%%%%%%%%%%%%%%%%%%%%%%%%%%%%%%%%%%%%%%%%%%%%%%
\section{Performance Summary}
\label{sec:string-performance}
%%%%%%%%%%%%%%%%%%%%%%%%%%%%%%%%%%%%%%%%%%%%%%%%%%%%%%%%%%%%%%%%%%%%%%%%%%%%%%%

\begin{table}[h]
\centering
\small
\begin{tabular}{lccc}
\textbf{Operation} & \textbf{Str} & \textbf{String} & \textbf{FixedString} \\
\hline
Create empty & --- & O(1) & O(1) \\
Create from C str & O(n) & O(n) & O(n) \\
Length & O(1) & O(1) & O(1) \\
Index byte & O(1) & O(1) & O(1) \\
Slice & O(1) & --- & --- \\
Push byte & --- & O(1)* & O(1) \\
Push string & --- & O(n)* & O(n) \\
Equals & O(n) & O(n) & O(n) \\
Find substring & O(nm) & O(nm) & O(nm) \\
\end{tabular}
\caption{String operation complexity. *Amortized for String (may reallocate).}
\end{table}

\begin{comingfromc}
Key differences from C strings:
\begin{itemize}
    \item \textbf{Length stored}: No need to scan for null terminator
    \item \textbf{Bounds checking}: Index operations return \texttt{Option}
    \item \textbf{RAII cleanup}: \texttt{String} and \texttt{BStr} free memory automatically
    \item \textbf{Safe slicing}: No pointer arithmetic errors
    \item \textbf{Type safety}: \texttt{Str} vs \texttt{String} distinguish borrowed vs owned
\end{itemize}
\end{comingfromc}
